\documentclass[a4paper]{arrowhead}

\usepackage[yyyymmdd]{datetime}
\usepackage{etoolbox}
\usepackage[T1]{fontenc}
\usepackage[utf8]{inputenc}
\usepackage{multirow}
\usepackage{longtable}
\usepackage{amssymb}
\usepackage{float}
\usepackage{tikz}
\usetikzlibrary{positioning,arrows.meta,fit,calc,shapes.misc,shapes.geometric}

\renewcommand{\dateseparator}{-}
\setlength{\emergencystretch}{3em}
\lstset{basicstyle=\ttfamily\footnotesize,breaklines=true,columns=fullflexible,frame=single,rulecolor=\color{gray!50},backgroundcolor=\color{gray!6}}
\tikzset{
  ahsys/.style={draw=ArrowheadBlue,very thick,fill=ArrowheadBlue!6,rounded corners=2pt,minimum width=3.2cm,minimum height=1.1cm,align=center,font=\small\bfseries},
  ahext/.style={draw=ArrowheadGrey,thick,dashed,fill=gray!5,rounded corners=2pt,minimum width=3.2cm,minimum height=1.0cm,align=center,font=\small},
  ahsrv/.style={draw=ArrowheadPurple,thick,fill=ArrowheadPurple!8,rounded corners=6pt,minimum height=0.55cm,inner xsep=4pt,align=center,font=\scriptsize\ttfamily},
  ahcsm/.style={draw=ArrowheadTeal,thick,fill=ArrowheadTeal!8,rounded corners=6pt,minimum height=0.55cm,inner xsep=4pt,align=center,font=\scriptsize\ttfamily},
  ahdb/.style={draw=ArrowheadOrange,thick,fill=ArrowheadOrange!8,cylinder,shape border rotate=90,aspect=0.25,minimum height=0.9cm,align=center,font=\scriptsize},
  ahflow/.style={-{Stealth[length=2.2mm]},thick,ArrowheadGrey},
  ahlabel/.style={font=\scriptsize,fill=white,inner sep=1pt,align=center},
}

%% Special references
\newcommand{\mref}[1]{{\textcolor{ArrowheadPurple}{\hyperref[sec:model:#1]{#1}}}}
\newcommand{\oref}[1]{{\textcolor{ArrowheadBlue}{\hyperref[sec:operations:#1]{#1}}}}
\newcommand{\pdef}[1]{{\textcolor{ArrowheadGrey}{#1\label{sec:model:primitives:#1}\label{sec:model:primitives:#1s}\label{sec:model:primitives:#1es}}}}
\newcommand{\pref}[1]{{\textcolor{ArrowheadGrey}{\hyperref[sec:model:primitives:#1]{#1}}}}

\newrobustcmd\osubsection[3]{
  \addtocounter{subsection}{1}
  \addcontentsline{toc}{subsection}{\protect\numberline{\thesubsection}operation \textcolor{ArrowheadBlue}{#1}}
  \subsection*{
    \thesubsection\quad
    operation
    \textcolor{ArrowheadBlue}{#1}
    (\notblank{#2}{\mref{#2}}{})
    \notblank{#3}{: \mref{#3}}{}
  }
  \label{sec:operations:#1}
}
\newrobustcmd\msubsection[2]{
  \addtocounter{subsection}{1}
  \addcontentsline{toc}{subsection}{\protect\numberline{\thesubsection}#1 \textcolor{ArrowheadPurple}{#2}}
  \subsection*{\thesubsection\quad#1 \textcolor{ArrowheadPurple}{#2}}
  \label{sec:model:#2} \label{sec:model:#2s} \label{sec:model:#2es}
}

\begin{document}

%% Arrowhead Document Properties
\ArrowheadTitle{workflowProduct}
\ArrowheadServiceID{workflowProduct}
\ArrowheadType{Service Description}
\ArrowheadTypeShort{SD}
\ArrowheadVersion{5.4.0}
\ArrowheadDate{2026-11-23}
\ArrowheadAuthor{Documentation maintenance}
\ArrowheadStatus{REVIEW DRAFT}
\ArrowheadContact{Repository maintainers}
\ArrowheadFooter{\href{www.arrowhead.eu}{www.arrowhead.eu}}
\ArrowheadSetup
%%

%% Front Page
\begin{center}
  \vspace*{1cm}
  \huge{\arrowtitle}

  \vspace*{0.2cm}
  \LARGE{\arrowtype}
  \vspace*{1cm}

  \Large{Service ID: \textit{``\arrowid''}}
  \vspace*{\fill}

  

  \vspace*{1cm}
  \vspace*{\fill}

  % Front Page Abstract
  \begin{abstract}
    This document describes the \texttt{workflowProduct} service of the Eclipse Arrowhead workflow support systems, which lets a WorkflowManager retrieve a product's recipe fragment and store the results back. It is provided by an Intelligent Product (consumed by the WorkflowManager), service version 1.0.0, documentation edition 5.4.0.
  \end{abstract}

  \vspace*{1cm}
\end{center}

\newpage
%%

%% Table of Contents
\tableofcontents
\newpage
%%

\section{Overview}
\label{sec:overview}
The \texttt{workflowProduct} service is \emph{consumed} by the WorkflowManager and \emph{provided}
by an Intelligent Product, its digital twin, or an adapter to the factory control system. It realises the principle
that data belongs at its source~\cite{derhamy2018}: the recipe and the as-built record belong to the product, the
WorkflowManager only works on a copy. This document specifies the contract a product must implement to be used with
\texttt{workflowInstance.start} and its \texttt{product} field.

The rest of this document is organized as follows.
In Section~\ref{sec:operations}, we describe the abstract message operations provided by the service.
In Section~\ref{sec:model}, we end the document by presenting the data types used by the mentioned operations.

\subsection{Significant Prior Art}
Derhamy's RetrieveProcessStep and StoreProcessStepResult interfaces of the Intelligent
Product~\cite{derhamy2018workflow}.

\subsection{How This Service Is Meant to Be Used}
When a session is started with \texttt{product} set, the WorkflowManager calls \oref{retrieve-process-step}
with its workstation identifier, validates and registers the returned recipe like any other recipe, and starts the
session with the returned parameters and subject. At Z-Close it calls \oref{store-process-step-result} with the stored
results; failures to store are logged in the worklog and do not change the session.

\subsection{Important Delimitations and Dependencies}
\label{sec:delimitations}
\begin{itemize}
\item Only the consumer side is implemented (\texttt{ProductGateway}); no product system exists in this workspace yet. The contract is a proposal (open decision 4 of the architecture, link to AAS/IEC~63278).
\item The product's acceptance or rejection of the result (\texttt{REJECTED}) is not yet acted upon.
\item Failed sessions are not stored back.
\end{itemize}

\paragraph{Authorization.} The product authorizes the WorkflowManager for both operations.

\newpage
\section{Service Interface}
\label{sec:operations}
This section describes the interface of the \texttt{workflowProduct} service. Each subsection names an abstract operation,
an input type and an output type, in that order. The input type is named inside parentheses, while the output type is
preceded by a colon. Input and output types are only denoted when accepted or returned, respectively, by the operation
in question. All abstract data types named in this section are defined in Section~\ref{sec:model}. The concrete
interface (HTTP/TLS/JSON) is specified in the IDD of the same name.

\begin{figure}[ht!]
  \centering
  \resizebox{0.92\textwidth}{!}{
\begin{tikzpicture}
\node[ahsys] (prov) at (4.6,0) {Intelligent Product\\[-1pt]{\normalfont\scriptsize $\ll$provider$\gg$}};
\node[ahsrv] (srv) at (1.6,0) {workflowProduct};
\draw[thick,ArrowheadPurple] (srv.east) -- (prov.west);
\node[ahext] (c0) at (-3.0,0.00) {WorkflowManager\\[-1pt]{\scriptsize $\ll$consumer$\gg$}};
\draw[ahflow] (c0.east) -- (srv.west);
\end{tikzpicture}
  }
  \caption{SysML block view of the \texttt{workflowProduct} service: provider and typical consumers.}
  \label{fig:SD}
\end{figure}

The following interface operations are available.

{\small
\begin{longtable}{|p{4.4cm}|p{5.6cm}|p{5.8cm}|}

\hline
\rowcolor{gray!33} Operation & Input & Output \\ \hline
\endhead
\oref{retrieve-process-step} & \mref{ProcessStepRequest} & \mref{ProcessOrder} \\ \hline
\oref{store-process-step-result} & \mref{ProcessStepResult} & -- \\ \hline
\end{longtable}
}

\osubsection{retrieve-process-step}{ProcessStepRequest}{ProcessOrder}
Returns the recipe fragment, parameters and subject for the requesting workstation.
\osubsection{store-process-step-result}{ProcessStepResult}{}
Stores the results of a completed session at the product.

\newpage
\section{Information Model}
\label{sec:model}
Here, all data objects that can be part of the \texttt{workflowProduct} service are listed in alphabetic order. Each
subsection, which describes one type of object, begins with the \textit{struct} keyword (a collection of named fields,
each with its own data type), the \textit{enum} keyword (a closed set of values) or the \textit{alias} keyword. As a
complement to the explicitly defined types in this section, there is also a list of implicit primitive types in
Section~\ref{sec:model:primitives}.

\msubsection{struct}{Action}
An action $\alpha$: a group of steps run in one work centre.

{\small
\begin{longtable}{|p{3.9cm}|p{4.4cm}|p{7.5cm}|}

\hline
\rowcolor{gray!33} Field & Type & Description \\ \hline
\endhead
\texttt{name} & \pref{Name} & Action name. \\ \hline
\texttt{next} & \pref{Name} & The action that becomes active when this one finishes; absent for the last action. Optional. \\ \hline
\texttt{workCentre} & \pref{String} & Informative work-centre identifier. Optional. \\ \hline
\texttt{dispatch} & \mref{DispatchMode} & Dispatch mode of this action's steps. Optional. \\ \hline
\texttt{steps} & \pref{List}$<$\mref{Step}$>$ & The steps $S_\alpha$. \\ \hline
\end{longtable}
}

\msubsection{enum}{DispatchMode}
How a step reaches an executor.

{\small
\begin{longtable}{|p{3.9cm}|p{11.9cm}|}

\hline
\rowcolor{gray!33} Value & Meaning \\ \hline
\endhead
\texttt{push} & The WorkflowManager selects an executor and calls \texttt{workflowExecution.start} (default). \\ \hline
\texttt{pull} & The WorkflowManager publishes the work item; an executor claims it with \texttt{workflowWorkItem.claim}. \\ \hline
\texttt{event} & Reserved for EventHandler notification; currently handled as \texttt{pull}. \\ \hline
\end{longtable}
}

\msubsection{struct}{ExecutorRequirement}
$\eta(s)$: constraints on the executor that runs a step.

{\small
\begin{longtable}{|p{3.9cm}|p{4.4cm}|p{7.5cm}|}

\hline
\rowcolor{gray!33} Field & Type & Description \\ \hline
\endhead
\texttt{capability} & \pref{String} & A capability the executor must declare in \texttt{workflowCapabilities}. Optional. \\ \hline
\texttt{workstation} & \pref{String} & The executor's workstation. Optional. \\ \hline
\texttt{system} & \pref{SystemName} & A named executor system. Optional. \\ \hline
\end{longtable}
}

\msubsection{enum}{LeaseMode}
Role of a step for a shared resource.

{\small
\begin{longtable}{|p{3.9cm}|p{11.9cm}|}

\hline
\rowcolor{gray!33} Value & Meaning \\ \hline
\endhead
\texttt{reserve} & Reserves the resource; its result carries the lease token. \\ \hline
\texttt{use} & Uses the resource; the executor renews the lease immediately before the invocation. \\ \hline
\texttt{release} & Releases the resource. \\ \hline
\end{longtable}
}

\msubsection{struct}{LeasePolicy}
The role of a step for a shared resource.

{\small
\begin{longtable}{|p{3.9cm}|p{4.4cm}|p{7.5cm}|}

\hline
\rowcolor{gray!33} Field & Type & Description \\ \hline
\endhead
\texttt{resource} & \pref{Name} & A declared resource. \\ \hline
\texttt{mode} & \mref{LeaseMode} & \texttt{reserve}, \texttt{use} or \texttt{release}. \\ \hline
\end{longtable}
}

\msubsection{enum}{OnError}
Reaction to an \texttt{ERROR} report.

{\small
\begin{longtable}{|p{3.9cm}|p{11.9cm}|}

\hline
\rowcolor{gray!33} Value & Meaning \\ \hline
\endhead
\texttt{rebind} & Exclude the executor for this session and bind the step again. \\ \hline
\texttt{abort} & Fail the session. \\ \hline
\end{longtable}
}

\msubsection{struct}{Parameter}
A declared session parameter.

{\small
\begin{longtable}{|p{3.9cm}|p{4.4cm}|p{7.5cm}|}

\hline
\rowcolor{gray!33} Field & Type & Description \\ \hline
\endhead
\texttt{name} & \pref{Name} & Parameter name, referenced as \texttt{\$\{param.name\}}. \\ \hline
\texttt{type} & \pref{String} & Informative type: \texttt{string}, \texttt{number}, \texttt{boolean} or \texttt{object}. Optional. \\ \hline
\texttt{required} & \pref{Boolean} & Whether \texttt{workflowInstance.start} refuses a session without it. Optional. \\ \hline
\end{longtable}
}

\msubsection{struct}{Policy}
Execution policy of a step. Every field other than \texttt{lease} and \texttt{localCloudOnly} is an extension outside the proven model.

{\small
\begin{longtable}{|p{3.9cm}|p{4.4cm}|p{7.5cm}|}

\hline
\rowcolor{gray!33} Field & Type & Description \\ \hline
\endhead
\texttt{quantity} & \pref{Integer} & Number of executions; only 1 is accepted (check M01). Optional. \\ \hline
\texttt{timeoutSeconds} & \pref{Duration} & Extension x-timeout: the step fails when it runs longer. Optional. \\ \hline
\texttt{nominalDurationSeconds} & \pref{Duration} & Expected duration, used to check lease durations (check L04). Optional. \\ \hline
\texttt{retry} & \mref{Retry} & Extension x-rebind: behaviour on \texttt{ERROR}. Optional. \\ \hline
\texttt{lease} & \mref{LeasePolicy} & The step's role for a shared resource. Optional. \\ \hline
\texttt{localCloudOnly} & \pref{Boolean} & Whether orchestration may return providers of other clouds (default true). Optional. \\ \hline
\end{longtable}
}

\msubsection{struct}{ProcessOrder}
The product's order for this workstation.

{\small
\begin{longtable}{|p{3.9cm}|p{4.4cm}|p{7.5cm}|}

\hline
\rowcolor{gray!33} Field & Type & Description \\ \hline
\endhead
\texttt{recipe} & \mref{Recipe} & The recipe fragment; validated and registered like any recipe. \\ \hline
\texttt{parameters} & \pref{JsonObject} & Session parameters $\rho$. Optional. \\ \hline
\texttt{subject} & \pref{String} & Product identity (subject of the session). Optional. \\ \hline
\end{longtable}
}

\msubsection{struct}{ProcessStepRequest}
Asks the product for its recipe fragment for one workstation.

{\small
\begin{longtable}{|p{3.9cm}|p{4.4cm}|p{7.5cm}|}

\hline
\rowcolor{gray!33} Field & Type & Description \\ \hline
\endhead
\texttt{workstationId} & \pref{String} & The \texttt{workflow.workstation.id} of the requesting WorkflowManager (empty when not configured). \\ \hline
\end{longtable}
}

\msubsection{struct}{ProcessStepResult}
What the WorkflowManager stores back to the product at Z-Close (data at its source).

{\small
\begin{longtable}{|p{3.9cm}|p{4.4cm}|p{7.5cm}|}

\hline
\rowcolor{gray!33} Field & Type & Description \\ \hline
\endhead
\texttt{workstationId} & \pref{String} & The workstation of the WorkflowManager. \\ \hline
\texttt{sessionId} & \pref{SessionId} & Session. \\ \hline
\texttt{recipe} & \pref{String} & \texttt{name@version}. \\ \hline
\texttt{status} & \mref{SessionStatus} & Final status; only \texttt{completed} sessions are stored back. \\ \hline
\texttt{results} & \pref{Map}$<$\pref{Name}, \pref{JsonValue}$>$ & $\xi$: the stored step results. \\ \hline
\end{longtable}
}

\msubsection{struct}{Recipe}
A workflow recipe $W$~\cite{varga2026}: named, versioned and immutable once registered. Actions run in sequence (\texttt{firstAction}, then each action's \texttt{next}); the steps of one action form a directed acyclic graph.

{\small
\begin{longtable}{|p{3.9cm}|p{4.4cm}|p{7.5cm}|}

\hline
\rowcolor{gray!33} Field & Type & Description \\ \hline
\endhead
\texttt{name} & \pref{Name} & Recipe name. \\ \hline
\texttt{version} & \pref{Version} & Recipe version; a registered name and version never changes content. \\ \hline
\texttt{description} & \pref{String} & Free text. Optional. \\ \hline
\texttt{firstAction} & \pref{Name} & The action $\alpha_0$ that is active when a session starts. \\ \hline
\texttt{parameters} & \pref{List}$<$\mref{Parameter}$>$ & Declared session parameters $\rho$. Optional. \\ \hline
\texttt{resources} & \pref{List}$<$\mref{Resource}$>$ & Shared resources (tools) used under leases. Optional. \\ \hline
\texttt{dispatch} & \mref{DispatchMode} & Default dispatch mode of the recipe; overridden per action or per session. Optional. \\ \hline
\texttt{actions} & \pref{List}$<$\mref{Action}$>$ & The actions $A$ with their steps. \\ \hline
\end{longtable}
}

\msubsection{struct}{Resource}
A shared resource that steps reserve, use and release under a lease (Proposition~1 of~\cite{varga2026}, lease variant).

{\small
\begin{longtable}{|p{3.9cm}|p{4.4cm}|p{7.5cm}|}

\hline
\rowcolor{gray!33} Field & Type & Description \\ \hline
\endhead
\texttt{name} & \pref{Name} & Resource name. \\ \hline
\texttt{leaseDurationSeconds} & \pref{Duration} & Lease duration granted by \texttt{workflowLease}; must exceed the longest tool-using step. Optional. \\ \hline
\texttt{provider} & \pref{SystemName} & Lease provider; the WorkflowManager when absent. Optional. \\ \hline
\end{longtable}
}

\msubsection{struct}{Retry}
Extension x-rebind.

{\small
\begin{longtable}{|p{3.9cm}|p{4.4cm}|p{7.5cm}|}

\hline
\rowcolor{gray!33} Field & Type & Description \\ \hline
\endhead
\texttt{maxAttempts} & \pref{Integer} & Maximum attempts of the step in one session. Optional. \\ \hline
\texttt{onError} & \mref{OnError} & What an \texttt{ERROR} report causes. Optional. \\ \hline
\end{longtable}
}

\msubsection{struct}{ServiceRequirement}
Requirements on the provider of the step's service, passed to orchestration.

{\small
\begin{longtable}{|p{3.9cm}|p{4.4cm}|p{7.5cm}|}

\hline
\rowcolor{gray!33} Field & Type & Description \\ \hline
\endhead
\texttt{serviceDefinition} & \pref{ServiceDefinition} & The service to consume. \\ \hline
\texttt{interfaces} & \pref{List}$<$\pref{String}$>$ & Accepted interface template names, e.g. \texttt{generic\_https}. Optional. \\ \hline
\texttt{minVersion} & \pref{String} & Lowest accepted service version. Optional. \\ \hline
\texttt{maxVersion} & \pref{String} & Highest accepted service version. Optional. \\ \hline
\texttt{metadata} & \pref{Map}$<$\pref{String}, \pref{JsonValue}$>$ & Metadata requirements on the service instance. Optional. \\ \hline
\texttt{qos} & \pref{Map}$<$\pref{String}, \pref{JsonValue}$>$ & QoS requirements forwarded to orchestration. Optional. \\ \hline
\texttt{requiresResults} & \pref{List}$<$\pref{Name}$>$ & Steps whose results the provider needs; each must be in $dep(s)$ (check D04). Optional. \\ \hline
\end{longtable}
}

\msubsection{enum}{SessionStatus}
Status of a session. \texttt{completed} and \texttt{failed} are final.

{\small
\begin{longtable}{|p{3.9cm}|p{11.9cm}|}

\hline
\rowcolor{gray!33} Value & Meaning \\ \hline
\endhead
\texttt{running} & Some rule may still apply. \\ \hline
\texttt{completed} & Z-Close fired: every action completed. \\ \hline
\texttt{failed} & S-FailBind, S-FailInvoke or an extension (abort, timeout) ended the session. \\ \hline
\end{longtable}
}

\msubsection{struct}{Step}
A step $s$: one service exchange.

{\small
\begin{longtable}{|p{3.9cm}|p{4.4cm}|p{7.5cm}|}

\hline
\rowcolor{gray!33} Field & Type & Description \\ \hline
\endhead
\texttt{name} & \pref{Name} & Step name, unique in the recipe. \\ \hline
\texttt{service} & \mref{ServiceRequirement} & $\theta(s)$: the service the step consumes. \\ \hline
\texttt{next} & \pref{List}$<$\pref{Name}$>$ & Successors in the same action (the transitions $s \rightarrow s'$). Optional. \\ \hline
\texttt{dep} & \pref{List}$<$\pref{Name}$>$ & $dep(s)$: steps whose results the step reads. Each must run before $s$. Optional. \\ \hline
\texttt{executor} & \mref{ExecutorRequirement} & $\eta(s)$: which executor may run the step; absent means any executor declaring the service. Optional. \\ \hline
\texttt{params} & \pref{Map}$<$\pref{String}, \pref{JsonValue}$>$ & Static parameters passed unchanged to the provider. Optional. \\ \hline
\texttt{input} & \pref{Template} & $\kappa_s$: the request body sent to the provider, composed from $\rho$ and $\xi|dep(s)$. Optional. \\ \hline
\texttt{policy} & \mref{Policy} & Execution policy (quantity, timeout, retry, lease, local cloud). Optional. \\ \hline
\texttt{guards} & \pref{Map}$<$\pref{Name}, \pref{Expression}$>$ & Conditions on the transitions to successors (extension x-guard). A successor without a guard is always taken. Optional. \\ \hline
\end{longtable}
}


\subsection{Primitives}
\label{sec:model:primitives}
Types and structures mentioned throughout this document that are assumed to be available to implementations of this
service. The concrete interpretation of each of these types must be provided by any IDD document claiming to implement
this service.

{\small
\begin{longtable}{|p{3.2cm}|p{12.6cm}|}

\hline
\rowcolor{gray!33} Type & Description \\ \hline
\endhead
\pdef{Boolean} & One out of \texttt{true} or \texttt{false}. \\ \hline
\pdef{Duration} & A length of time in seconds. \\ \hline
\pdef{Expression} & A side-effect free condition over the session parameters and stored results (Section on guard expressions). \\ \hline
\pdef{Integer} & A whole number. \\ \hline
\pdef{JsonObject} & A set of named values. \\ \hline
\pdef{JsonValue} & Any value: object, array, string, number, Boolean or null. \\ \hline
\pdef{List} & An ordered collection of items of the type given in angle brackets. \\ \hline
\pdef{Map} & A collection of key/value pairs of the types given in angle brackets. \\ \hline
\pdef{Name} & An identifier of a recipe, action, step, parameter or resource: a letter followed by at most 62 letters, digits, \texttt{-} or \texttt{\_}. \\ \hline
\pdef{ServiceDefinition} & The name of an Arrowhead service definition. Names are compared case-insensitively and without \texttt{-}, \texttt{\_} and spaces, so \texttt{temperature-read} and \texttt{temperatureRead} denote the same service. \\ \hline
\pdef{SessionId} & The identifier the WorkflowManager gives a session. \\ \hline
\pdef{String} & A chain of characters. \\ \hline
\pdef{SystemName} & The name of an Arrowhead system as it appears in its certificate and in the ServiceRegistry. \\ \hline
\pdef{Template} & A value in which strings of the form \texttt{\$\{param.$p$\}} or \texttt{\$\{result.$s$.path\}} are replaced by a session parameter or by (a part of) a stored step result. \\ \hline
\pdef{Version} & A recipe version in the form MAJOR.MINOR.PATCH. \\ \hline
\end{longtable}
}


\newpage
\bibliographystyle{IEEEtran}
\bibliography{bibliography}

\newpage

\section{Revision History}
\subsection{Amendments}

\noindent\begin{tabularx}{\textwidth}{| p{1cm} | p{2.2cm} | p{1.6cm} | X | p{3.6cm} |} \hline
\rowcolor{gray!33} No. & Date & Version & Subject of Amendments & Author \\ \hline
1 & 2026-11-23 & \arrowversion & Initial edition: service description of workflowProduct, written from the Java reference implementation (ah-Workflow-java). & Documentation maintenance \\ \hline
\end{tabularx}

\subsection{Quality Assurance}

\noindent\begin{tabularx}{\textwidth}{| p{1cm} | p{2.2cm} | p{1.6cm} | X |} \hline
\rowcolor{gray!33} No. & Date & Version & Approved by \\ \hline
1 & & \arrowversion & Pending review \\ \hline
\end{tabularx}

\end{document}
